\section{Methodology}


We describe the proposed GAD-MambaUNet by first introducing its core building block and then integrating it into an asymmetric U-shaped segmentation architecture. 
As shown in Fig.~\ref{fig:overall-network}, the left part presents the overall student network and the training-time DINOv3 teacher, while the right part details the proposed block-level design and the internal structure of Direction-Group Graph Selective Scan (DG-GSS). 
We first present DG-GSS, which enables structured interaction among scan directions and channel groups. Then, we describe how DG-GSS is embedded into the overall GAD-MambaUNet architecture. Finally, we introduce the training-time DINOv3 supervision with Gradient-Adaptive Distillation (GAD), where the teacher model is used only during training and removed during inference.


\subsection{Direction-Group Graph Selective Scan}

As illustrated in the lower-right part of Fig.~\ref{fig:overall-network}, DG-GSS is designed to enhance grouped multi-directional selective scanning by introducing graph-based interaction among direction--group responses. Grouped multi-directional selective scanning reduces the cost of Vision Mamba by dividing feature channels into multiple groups and processing them along several traversal directions. However, in conventional grouped scanning, each direction--group branch is processed independently and fused only after selective scanning. This independent processing limits the exchange of complementary information across scan directions and channel subspaces. To address this issue, DG-GSS regards each scan-direction and channel-group response as a graph node and performs structured message passing before the final multi-directional fusion.

Let $\mathbf{X}\in\mathbb{R}^{B\times D_i\times h\times w}$ denote the inner state-space feature of DG-GSS, where $D_i$ is the inner channel dimension. We adopt four Cross2D traversal directions and partition the inner channels into $M$ groups. With $d=D_i/M$ and $L=hw$, cross scanning produces
\begin{equation}
\mathbf{X}_{\mathrm{scan}}\in\mathbb{R}^{B\times K\times M\times d\times L},\qquad K=4.
\end{equation}

For each direction--group pair, the selective-scan parameters $(\Delta,\mathbf{B},\mathbf{C})$ are generated by independent grouped projections, and selective scanning is applied along the corresponding sequence. The directional outputs are then realigned to the original spatial coordinate system:
\begin{equation}
\mathbf{Y}\in\mathbb{R}^{B\times K\times M\times d\times h\times w}.
\end{equation}

Each tensor $\mathbf{Y}_{k,m}$ corresponds to the response of scan direction $k$ and channel group $m$. We regard every direction--group response as a graph node and obtain its node descriptor by global average pooling:
\begin{equation}
\mathbf{r}_{k,m}
=
\frac{1}{hw}
\sum_{u=1}^{h}
\sum_{v=1}^{w}
\mathbf{Y}_{k,m,:,u,v}.
\end{equation}
After layer normalization and linear projection, the node descriptor is transformed into $\mathbf{h}_{k,m}$.

In our implementation, we adopt a factorized direction-group graph, where two nodes are connected if they share the same scan direction or the same channel group, and self-connections are included. For two connected nodes $i$ and $j$, additive graph-attention logits~\cite{velickovic2018gat} are computed as
\begin{equation}
e_{ij}
=
\operatorname{LeakyReLU}
\left(
\mathbf{a}_{s}^{\top}\mathbf{h}_{i}
+
\mathbf{a}_{t}^{\top}\mathbf{h}_{j}
\right).
\end{equation}
The attention coefficients are obtained by masked softmax over the neighborhood:
\begin{equation}
\alpha_{ij}
=
\frac{\exp(e_{ij})}
{\sum_{j'\in\mathcal{N}(i)}\exp(e_{ij'})},
\qquad j\in\mathcal{N}(i).
\end{equation}

Graph-based message passing is then performed on the full spatial feature maps:
\begin{equation}
\widetilde{\mathbf{Y}}_i
=
\mathbf{Y}_i
+
\gamma
\sum_{j\in\mathcal{N}(i)}
\alpha_{ij}V(\mathbf{Y}_j),
\end{equation}
where $V(\cdot)$ is a shared $1\times1$ value projection and $\gamma$ is a learnable scaling parameter initialized to zero. This initialization makes DG-GSS start from the original grouped selective-scan behavior and progressively learn direction--group information exchange during training.

Finally, the enhanced group responses are concatenated along the channel dimension for each direction, and the four aligned directional feature maps are summed:
\begin{equation}
\mathbf{Z}
=
\sum_{k=1}^{K}
\operatorname{Concat}_{m=1}^{M}
\left(
\widetilde{\mathbf{Y}}_{k,m}
\right).
\end{equation}
Output normalization is then applied to obtain the DG-GSS output.
In this way, DG-GSS preserves the efficiency of grouped multi-directional selective scanning while explicitly modeling the structural relationship between scan directions and channel groups.



\subsection{GAD-MambaUNet Architecture}

After defining DG-GSS, we embed it into a lightweight block and integrate the block into an asymmetric U-shaped segmentation network. As shown in the upper-right part of Fig.~\ref{fig:overall-network}, the DG-GSS block consists of three residual components: a pre-scan local mixing stage, the proposed DG-GSS layer, and a post-scan local refinement stage. Given an input feature $\mathbf{x}$, the block is formulated as
\begin{equation}
\begin{aligned}
\mathbf{u}_1 &= \mathbf{x}+\operatorname{FFN}_0(\operatorname{DW}_0(\mathbf{x})), \\
\mathbf{u}_2 &= \mathbf{u}_1+\operatorname{DG\text{-}GSS}(\mathbf{u}_1), \\
\mathbf{y} &= \mathbf{u}_2+\operatorname{FFN}_1(\operatorname{DW}_1(\mathbf{u}_2)),
\end{aligned}
\end{equation}
where $\operatorname{DW}(\cdot)$ denotes depth-wise convolution for local mixing and $\operatorname{FFN}(\cdot)$ denotes a lightweight feed-forward refinement module. The local operators enhance short-range texture and channel interactions, while DG-GSS provides efficient long-range context aggregation through direction--group graph interaction.

The overall GAD-MambaUNet follows a five-level encoder--decoder architecture, as shown in the left part of Fig.~\ref{fig:overall-network}. Given an input image $\mathbf{x}\in\mathbb{R}^{3\times H\times W}$, the network predicts a binary segmentation logit map $\mathbf{z}\in\mathbb{R}^{1\times H\times W}$. In the base configuration, the channel widths are set to $[16,32,64,96,160]$. To balance local boundary preservation and global context modeling, the first three encoder stages and the last three decoder stages use lightweight multi-kernel inverted residual (MKIR) blocks inherited from MK-UNet~\cite{rahman2025mkunet}, while the two deepest encoder stages and the first two decoder stages use DG-GSS blocks. Resolution and channel changes are performed outside these blocks through depth-wise downsampling with point-wise projection or point-wise projection followed by bilinear upsampling, keeping the feature width fixed inside each block.

For decoder reconstruction, skip connections transfer high-resolution encoder features to the corresponding decoder stages. Before fusion, a grouped attention gate filters the skip feature according to the decoder context, reducing irrelevant background responses. Let $\mathbf{g}$ denote the upsampled decoder feature and $\mathbf{s}$ the encoder feature at the same resolution. The gated skip feature is computed as
\begin{equation}
\boldsymbol{\alpha}=\sigma!\left(\psi!\left(\phi(W_g\mathbf{g}+W_s\mathbf{s})\right)\right),\qquad
\hat{\mathbf{s}}=\boldsymbol{\alpha}\odot\mathbf{s},
\end{equation}
where $W_g$ and $W_s$ are grouped convolutions, $\phi$ is ReLU, and $\psi$ maps the joint response to a spatial gate. The filtered skip feature $\hat{\mathbf{s}}$ is then fused with the decoder feature. The first decoder block operates at the deepest decoder level and provides the aligned student feature for the training-time DINOv3 supervision described in the next subsection.



\subsection{Training-Time DINOv3 Supervision with GAD}

Although DG-GSS improves contextual modeling inside the compact student network, training a lightweight segmentation model on limited medical datasets may still lead to insufficient semantic abstraction. Inspired by the DINOv3-based gradient-adaptive distillation strategy in RT-DETRv4~\cite{liao2025rtdetrv4}, we adapt training-time foundation-model supervision to lightweight medical image segmentation. A frozen DINOv3 ViT-B/16 model~\cite{simeoni2025dinov3} is used as the semantic teacher only during training, and the teacher branch is removed during inference.

As described in the previous subsection, the student feature used for distillation is extracted after the first decoder block and before the first upsampling operation. This feature is semantically rich, spatially compact, and computationally inexpensive for feature alignment. Let $\mathbf{F}_{s}$ denote the projected student feature, and let $\mathbf{F}_{t}$ denote the corresponding DINOv3 teacher feature. The student feature is projected to the teacher dimension by a $1\times1$ convolution followed by batch normalization when their channel dimensions are different. If necessary, the teacher feature is bilinearly resized to match the spatial resolution of $\mathbf{F}_{s}$.

After spatial alignment, both feature maps are flattened into $N$ spatial tokens and $\ell_2$-normalized along the channel dimension. We use a cosine feature distillation loss:
\begin{equation}
\mathcal{L}_{\mathrm{dist}}
=
\frac{1}{N}\sum_{n=1}^{N}
\left[
1-
\frac{
\mathbf{f}_{s,n}^{\top}\mathbf{f}_{t,n}
}{
\left\|\mathbf{f}_{s,n}\right\|_2
\left\|\mathbf{f}_{t,n}\right\|_2
+\epsilon
}
\right],
\end{equation}
where $\mathbf{f}_{s,n}$ and $\mathbf{f}_{t,n}$ denote the $n$-th student and teacher tokens, respectively, and $\epsilon$ is a small constant for numerical stability.

The segmentation objective follows the structure-aware weighted BCE and weighted IoU loss~\cite{fan2020pranet}:
\begin{equation}
\mathcal{L}_{\mathrm{seg}}
=
\mathcal{L}_{\mathrm{wbce}}
+
\mathcal{L}_{\mathrm{wiou}}.
\end{equation}
The overall training objective is
\begin{equation}
\mathcal{L}
=
\mathcal{L}_{\mathrm{seg}}
+
\lambda_e\mathcal{L}_{\mathrm{dist}},
\end{equation}
where $\lambda_e$ is the distillation coefficient at epoch $e$.

Using a fixed distillation coefficient assumes that the proper teacher supervision strength remains unchanged throughout training. However, the optimization state of the student changes over time, and the gradient contribution of the teacher-supervised decoder block may become either too weak or too dominant. To balance semantic supervision and the primary segmentation objective, GAD dynamically updates $\lambda_e$ according to the gradient share of the aligned decoder block. After back-propagation and before gradient clipping, we compute
\begin{equation}
q_e
=
100\times
\frac{
\sum_{\theta_i\in\Theta_a}
\left\|\nabla_{\theta_i}\mathcal{L}\right\|_1
}{
\sum_{\theta_j\in\Theta}
\left\|\nabla_{\theta_j}\mathcal{L}\right\|_1
+\epsilon
},
\end{equation}
where $\Theta_a$ denotes the parameters of the aligned decoder block and $\Theta$ denotes all trainable student parameters.

Given a target gradient-share interval $[\rho-\delta,\rho+\delta]$, GAD adjusts the next-epoch distillation coefficient only when $q_e$ falls outside this interval. Let $q^{*}$ be the nearest target boundary, and define $p_e=q_e/100$ and $p^{*}=q^{*}/100$. The odds-ratio adjustment is computed as
\begin{equation}
r_e
=
\frac{
p^{*}(1-p_e)
}{
p_e(1-p^{*})+\epsilon
}.
\end{equation}
The distillation coefficient for the next epoch is updated by
\begin{equation}
\lambda_{e+1}
=
\operatorname{clip}
\left(
\lambda_e
\operatorname{clip}(r_e,r_{\min},r_{\max}),
\lambda_{\min},
\lambda_{\max}
\right),
\end{equation}
where $r_{\min}$ and $r_{\max}$ bound the per-epoch adjustment, and $\lambda_{\min}$ and $\lambda_{\max}$ constrain the overall distillation strength.

Training begins with a segmentation-only burn-in stage, followed by a linear distillation warm-up. GAD is activated after warm-up to adaptively regulate the teacher contribution. During inference, the DINOv3 teacher, the feature projection head, and the GAD update are all discarded. Therefore, the proposed training-time supervision improves the compact student without increasing deployment cost.


% \section{Methodology}

% \subsection{Overview}

% Fig.~\ref{fig:overall-network} presents GAD-MambaUNet. Given an image $\mathbf{x}\in\mathbb{R}^{3\times H\times W}$, the student produces a binary logit map $\mathbf{z}\in\mathbb{R}^{1\times H\times W}$. The base model uses five resolution levels with channel widths $[16,32,64,96,160]$. The first three encoder blocks and the last three decoder blocks are multi-kernel inverted residual (MKIR) blocks. The two deepest encoder blocks and the first two decoder blocks are grouped Mamba blocks. Resolution and channel changes are performed outside these blocks by depth-wise downsampling followed by point-wise projection, or by point-wise projection followed by bilinear upsampling. This separation keeps the state-space width fixed within each block.

% At every decoder level, channel and spatial attention refine the decoder feature, while a grouped attention gate filters the corresponding skip connection before additive fusion. The first decoder block operates at $H/32\times W/32$ and supplies the student feature used for DINOv3 supervision. This location is semantically rich, inexpensive, and directly aligned with the block monitored by GAD.

% \subsection{Asymmetric Local--Global Student}

% At high resolution, the dominant requirement is local boundary reconstruction. We therefore retain the MKIR design of MK-UNet~\cite{rahman2025mkunet}: a $1\times1$ expansion, parallel depth-wise convolutions with kernels $\mathcal{K}=\{1,3,5\}$, summation, channel shuffle, and a $1\times1$ projection. For feature $\mathbf{x}$, the block is written compactly as
% \begin{equation}
% \operatorname{MKIR}(\mathbf{x})=
% P_2\!\left(\operatorname{Shuffle}\!\left[
% \sum_{k\in\mathcal{K}}D_k\!\left(P_1(\mathbf{x})\right)\right]\right),
% \end{equation}
% where $P_1$ and $P_2$ denote point-wise expansion and projection and $D_k$ is a depth-wise convolution with kernel size $k$. The block supplies multi-receptive-field local evidence at low cost.

% At low resolution, the receptive field can be enlarged more efficiently. A grouped Mamba block first applies depth-wise local mixing and a grouped shuffle feed-forward network (FFN), then DG-GSS, and finally a second local mixing--FFN pair. Each component is residual:
% \begin{equation}
% \begin{aligned}
% \mathbf{u}_1 &= \mathbf{x}+\operatorname{FFN}_0(\operatorname{DW}_0(\mathbf{x})),\\
% \mathbf{u}_2 &= \mathbf{u}_1+\operatorname{DG\text{-}GSS}(\mathbf{u}_1),\\
% \mathbf{y} &= \mathbf{u}_2+\operatorname{FFN}_1(\operatorname{DW}_1(\mathbf{u}_2)).
% \end{aligned}
% \end{equation}
% The local operators therefore refine state-space features instead of competing with them.

% For skip fusion, let $\mathbf{g}$ denote the upsampled decoder feature and $\mathbf{s}$ the encoder feature at the same resolution. A grouped attention gate computes
% \begin{equation}
% \boldsymbol{\alpha}=\sigma\!\left(\psi\!\left(\phi(W_g\mathbf{g}+W_s\mathbf{s})\right)\right),\qquad
% \hat{\mathbf{s}}=\boldsymbol{\alpha}\odot\mathbf{s},
% \end{equation}
% where $W_g$ and $W_s$ are grouped convolutions, $\phi$ is ReLU, and $\psi$ maps the joint response to a spatial gate. The decoder feature and $\hat{\mathbf{s}}$ are then added.

% \subsection{Direction-Group Graph Selective Scan}

% Let $\mathbf{X}\in\mathbb{R}^{B\times D\times h\times w}$ be the input to DG-GSS. We use four Cross2D traversal directions and partition the inner state-space channels into $M$ groups. With $d=D/M$ and $L=hw$, cross scanning yields
% \begin{equation}
% \mathbf{X}_{\mathrm{scan}}\in\mathbb{R}^{B\times K\times M\times d\times L},\qquad K=4.
% \end{equation}
% Every direction--group pair independently produces its selective-scan parameters $(\Delta,\mathbf{B},\mathbf{C})$. Selective scan is applied to each pair and the four outputs are spatially realigned to the original coordinate system:
% \begin{equation}
% \mathbf{Y}\in\mathbb{R}^{B\times K\times M\times d\times h\times w}.
% \end{equation}

% The tensor $\mathbf{Y}_{k,m}$ is a response from direction $k$ and channel group $m$. We regard it as a graph node and obtain a node descriptor by global average pooling:
% \begin{equation}
% \mathbf{r}_{k,m}=\frac{1}{hw}\sum_{u,v}\mathbf{Y}_{k,m,:,u,v}.
% \end{equation}
% After layer normalization and linear projection, additive graph-attention logits~\cite{velickovic2018gat} are
% \begin{equation}
% e_{ij}=\operatorname{LeakyReLU}\left(\mathbf{a}_{s}^{\top}\mathbf{h}_i+
% \mathbf{a}_{t}^{\top}\mathbf{h}_j\right).
% \end{equation}
% Our factorized adjacency connects two nodes when they have the same scan direction or the same channel-group identity, including self-connections. Thus, a node receives both cross-group and cross-direction evidence without constructing an unconstrained fully connected graph. With masked softmax weights $a_{ij}$ and shared $1\times1$ value projection $V(\cdot)$, graph propagation is
% \begin{equation}
% \widetilde{\mathbf{Y}}_i=\mathbf{Y}_i+\gamma\sum_{j\in\mathcal{N}(i)}a_{ij}V(\mathbf{Y}_j),
% \end{equation}
% where $\gamma$ is learnable and initialized to zero. DG-GSS consequently starts from the original grouped scan behavior and learns graph exchange progressively. Finally, groups are concatenated and the four aligned directions are summed before output normalization.

% \subsection{DINOv3 Supervision and Gradient-Adaptive Distillation}

% A frozen DINOv3 ViT-B/16 teacher provides dense features only during training. The feature $\mathbf{F}_s$ after the first decoder block is projected by a $1\times1$ convolution and batch normalization to the teacher dimension. The teacher feature $\mathbf{F}_t$ is bilinearly resized when necessary. After flattening each map into spatial tokens and applying $\ell_2$ normalization, the feature loss is
% \begin{equation}
% \mathcal{L}_{\mathrm{dist}}=
% \frac{1}{N}\sum_{n=1}^{N}\left[1-
% \frac{\mathbf{f}_{s,n}^{\top}\mathbf{f}_{t,n}}
% {\|\mathbf{f}_{s,n}\|_2\|\mathbf{f}_{t,n}\|_2}\right].
% \end{equation}
% The segmentation objective follows the structure-aware weighted BCE and weighted IoU loss~\cite{fan2020pranet}:
% \begin{equation}
% \mathcal{L}=\mathcal{L}_{\mathrm{seg}}+\lambda_e\mathcal{L}_{\mathrm{dist}},
% \qquad
% \mathcal{L}_{\mathrm{seg}}=\mathcal{L}_{\mathrm{wbce}}+\mathcal{L}_{\mathrm{wiou}}.
% \end{equation}

% Using a fixed $\lambda_e$ assumes that the appropriate teacher pressure remains unchanged throughout training. Instead, after backpropagation and before gradient clipping, GAD measures the gradient share of the aligned decoder block:
% \begin{equation}
% q_e=100\times
% \frac{\sum_{\theta_i\in\Theta_a}\|\nabla_{\theta_i}\mathcal{L}\|_1}
% {\sum_{\theta_j\in\Theta}\|\nabla_{\theta_j}\mathcal{L}\|_1}.
% \end{equation}
% Here, $\Theta_a$ is the parameter set of decoder1 and $\Theta$ contains all student parameters. If $q_e$ falls outside the target interval $[\rho-\delta,\rho+\delta]$, GAD selects the opposite interval boundary as $q^*$ and computes an odds-ratio adjustment:
% \begin{equation}
% r_e=\frac{p^{*}(1-p_e)}{p_e(1-p^{*})},
% \quad p_e=q_e/100,\quad p^{*}=q^{*}/100.
% \end{equation}
% The next base distillation coefficient is
% \begin{equation}
% \lambda_{e+1}=\operatorname{clip}\left(
% \lambda_e\operatorname{clip}(r_e,r_{\min},r_{\max}),
% \lambda_{\min},\lambda_{\max}\right).
% \end{equation}
% The multiplicative adjustment is bounded per epoch and the coefficient is clipped to $[\lambda_{\min},\lambda_{\max}]$ for stability. Training begins with segmentation-only burn-in, follows with a linear distillation warm-up, and activates GAD after warm-up. At inference, the teacher, feature projection head, and GAD update are discarded.
